\documentclass[11pt, a4paper]{article}

% --- PACCHETTI ---
\usepackage[utf8]{inputenc}
\usepackage[italian]{babel}
\usepackage[margin=2.5cm]{geometry} % Imposta i margini della pagina
\usepackage{graphicx}               % Per inserire il logo
\usepackage{tabularx}               % Per tabelle larghe quanto la pagina
\usepackage[table]{xcolor}          % Per colorare le celle delle tabelle
\usepackage[hidelinks]{hyperref}    % Per stile link cliccabili
\usepackage{lmodern}      % Rende i font scalabili all'infinito
\usepackage{anyfontsize}  % Permette di usare dimensioni personalizzate
\usepackage{fancyhdr}

% Aggiungi la linea di separazione sotto l'intestazione
\renewcommand{\headrulewidth}{0.4pt}

\begin{document}

% ==============================================================================
% FRONTESPIZIO (PAGINA 1)
% ==============================================================================
\begin{titlepage}
    \centering
    \vspace*{0.5cm}

    % LOGHI
    \includegraphics[width=0.45\textwidth]{../../../assets/logo_unipd.jpg}
    \hfill
    \includegraphics[width=0.45\textwidth]{../../../assets/logo.png}\par
    \vspace{0.5cm}

    % NOME DEL GRUPPO
    {\fontsize{45}{50}\sffamily\selectfont \textbf{On-By-One}\par} 
    \vspace{1cm}

    % TITOLO DEL VERBALE
    {\huge \bfseries Lettera di presentazione \par}
    
    \vspace{1.5cm}

    % MEMBRI DEL GRUPPO
    {\large
    \textbf{On-By-One} (Gruppo 1):\par
    \vspace{0.3cm}
    \begin{tabular}{l l}
        Sara Ballardin & 2101048 \\
        Shuhui Chen & 2138001 \\
        Simone Giuliano & 2145591 \\
        Alvise Greggio & 2141004 \\
        Daniele Nuzzaci & 2139039 \\
        Federico Temporin & 2146764 \\
        Tommaso Zorzetto & 2148629 \\
    \end{tabular}
    \par}

    \vfill

    % TABELLA INFORMAZIONI DOCUMENTO E DESTINATARI
    \begin{table}[h]
        \centering
        \arrayrulecolor{gray!50}
        \renewcommand{\arraystretch}{1.5}
        \begin{tabularx}{\textwidth}{|c|X|}
            \hline
                \rowcolor{gray!10} \multicolumn{2}{|c|}{\textbf{Informazioni documento}} \\
            \hline
                \rowcolor{gray!10} \textbf{Data} & \multicolumn{1}{c|}{\textbf{Destinatari}} \\
            \hline
                2026-10-30 & \centering\arraybackslash On-By-One, Prof. Tullio Vardanega, Prof. Riccardo Cardin \\
            \hline
        \end{tabularx}
        \arrayrulecolor{black}
    \end{table}
    
    \vspace{0.5cm}
    
    % CONTATTI
    {\large Contatti: \href{mailto:onbyone.swe@gmail.com}{onbyone.swe@gmail.com} }
    
\end{titlepage}

\newpage

\section{Scopo documento} Con il presente documento, il team \textbf{On-By-One} intende formalizzare la propria candidatura per lo sviluppo del capitolato \textbf{C1 AURA} (AI for Understanding, Reasoning and Automation), proposto da \textbf{Vimar S.P.A.}. 
La documentazione a supporto della candidatura, comprendente la valutazione dei capitolati, la dichiarazione degli impegni e i verbali degli incontri esplorativi con le aziende proponenti, è consultabile presso il \href{https://on-by-one-team.github.io/On-By-One.github.io/}{sito web del team}\footnote{https://on-by-one-team.github.io/On-By-One.github.io/}. \\

\section{Motivazione scelta} La decisione di candidarci per questo progetto è maturata a seguito di un'attenta valutazione delle proposte disponibili, documentata nell'apposita \href{https://github.com/On-By-One-Team/On-By-One.github.io/blob/main/docs/pdf/candidatura/valutazione-capitolati.pdf}{analisi comparativa}\footnote{https://github.com/On-By-One-Team/On-By-One.github.io/blob/main/docs/pdf/candidatura/valutazione-capitolati.pdf}.

\noindent AURA si è distinto per la possibilità di coniugare lo sviluppo software con una componente di sperimentazione hardware, grazie all'impiego di un kit per impianti domotici smart e di moduli dedicati al testing e all'interazione fisica.

\noindent Questa caratteristica ci permetterà di mettere in pratica i principi dell'ingegneria del software affrontando l'integrazione tra intelligenza artificiale e standard domotici avanzati, per poi validare l'architettura direttamente sul campo.

\noindent Abbiamo inoltre valutato positivamente la complessità tecnica del progetto, incentrato su ambienti \textit{cloud} e servizi intelligenti distribuiti, che riteniamo coerente con le competenze del gruppo e stimolante per la crescita professionale di ciascun componente.

\noindent Un ulteriore elemento di valutazione è stato il rigore con cui la proposta è stata formulata: la completezza della documentazione, la chiarezza dei requisiti e la disponibilità dimostrata da Vimar S.P.A. durante gli incontri esplorativi costituiscono presupposti favorevoli per una collaborazione efficace e per una gestione consapevole delle attività progettuali.

\noindent Infine, la presenza di requisiti opzionali e di possibili evoluzioni future, quali la gestione multi-impianto e l'analisi dello storico dei dati, offre ulteriori opportunità di approfondimento tecnico, consentendo di valutare eventuali estensioni del sistema in funzione delle priorità e delle risorse disponibili. \\

\section{Conclusioni} Il preventivo dei costi finali, riportato nel documento di \href{https://github.com/On-By-One-Team/On-By-One.github.io/blob/main/docs/pdf/candidatura/dichiarazione-degli-impegni.pdf}{dichiarazione degli impegni}\footnote{https://github.com/On-By-One-Team/On-By-One.github.io/blob/main/docs/pdf/candidatura/dichiarazione-degli-impegni.pdf}, è stato calcolato secondo il regolamento previsto e ammonta a \textbf{€ 12.840,00}. La scadenza ultima di consegna è fissata al \textbf{19 marzo 2027}. \\

\noindent Alla luce delle considerazioni esposte, confermiamo il nostro interesse per il capitolato AURA e la volontà di affrontarlo con serietà, rigore metodologico e attenzione alla qualità del prodotto finale, valorizzando al contempo il confronto con l'azienda proponente. \\

\noindent Cordiali saluti, 

\noindent Il team On-By-One.
\end{document}